% sim_robustness.tex -- Results: robustness to wrong parameters (draft 4 Oct 2026).
% Sources: CURRENT_STATE 4.7 (E10, R0215-R0222, 10 Sep, attach demo, 0.6 kg ring), 4.8 (measure_rod_len,
% SIL R0308-R0327, 23 Sep, quadratic plant, notes 2026-09-23_geometry_autocal), 4.10 (kt_trim,
% R0405-R0432, notes 2026-09-23_kt_trim), decisions 2026-09-25 (height integral: R0557 0.661 vs
% R0558 0.600; pair 0.658 x2 vs 0.600 x3).
% All of these were flown with the tracker's node-0 box at 0.025 (default until 4 Oct), which is the
% +3.4 cm of the 0.634 baselines; the differences between arms are the result.
% T9 (tab:robust): R1038-R1067 (5 Oct, box 0, land fix), no-error baseline z_ki 0 R1013/R1023;
% generator tools/thesis_sim_tables.py table_robust.

\subsection{Robustness to wrong parameters}
\label{sec:sim_robustness}

The controller is given the payload's mass, the rod lengths and the attachment radius, and each UAV's
thrust map. On the rig each of these is measured, and each can be wrong. This subsection collects the
tests of what a wrong value does.

\paragraph{The current controller.} Table~\ref{tab:robust} gives the rig twin's hover at 1.0\,m, three
UAVs from a floor start, with one wrong value in the controller while the simulated rig keeps the true
ones. Each case was flown without and with the height integral.
\begin{itemize}
    \item A wrong payload mass moves the hover by about 1.25\,cm per percent of mass, the same in
    both directions. The integral removes it while the miss is inside its 0.25\,m gate, with a target
    shift that mirrors the error. A 20\,\% overestimate puts the ring 2\,mm over the gate, so the
    integral never starts and the ring holds 26\,cm high: a miss beyond the gate is not taken for a
    static offset. The bound matters as well: the 20\,\% underestimate needs 23\,cm of shift, which
    the earlier bound of 0.15\,m would have cut off, and the planner warns that the integral is near
    its bound.
    \item Measuring the rods at the hand-over absorbs a 5\,cm rod error completely. Without the
    measurement the ring sits 8\,cm high, and the integral does not help: a rod typed 5\,cm long
    measures as only partly taut, and the integral waits for every rod to be taut. For rods typed
    long, the measurement is therefore what makes the integral available; a rod typed short reads as
    taut, and that case was not flown.
    \item A 2\,cm error in the attachment radius has no visible effect.
\end{itemize}
The ring's tilt stayed within 1.0--1.3$^\circ$, slightly lower with the mass typed high, and all 28
flights landed.

\begin{table}[htbp]
\centering
\caption[Hover with a wrong model]{Hover at 1.0\,m with one wrong value in the controller's model,
Gazebo rig twin, three UAVs, $\epsilon = 0$, floor start. Each cell is the mean of two runs. Target
shift: the integral's correction to the height target at LAND. Tilt: mean over the four runs.
Landed: flights that landed without a tip. $^a$~The miss is over the integral's 0.25\,m gate, so the
integral never starts. $^b$~The unmeasured rods read as partly slack, so the integral's tautness
gate stays closed.}
\label{tab:robust}
% generated by tools/thesis_sim_tables.py
\begin{tabularx}{\textwidth}{>{\raggedright\arraybackslash}Xccccc}
\toprule
 & \multicolumn{2}{c}{Height error (cm)} & Target & & \\
\cmidrule(lr){2-3}
Error in the controller's model & no integral & integral & shift (cm) & Tilt (\si{\degree}) & Landed \\
\midrule
none & $+$0.9 & 0.0 & $-$0.8 & 1.2 & 4/4 \\
payload mass typed $+10$\,\% & $+$13.5 & 0.0 & $-$13.4 & 1.1 & 4/4 \\
payload mass typed $+20$\,\% & $+$25.8 & $+$25.8$^a$ & 0.0 & 1.0 & 4/4 \\
payload mass typed $-20$\,\% & $-$23.5 & 0.0 & $+$23.3 & 1.3 & 4/4 \\
rods typed $+5$\,cm, measured at the hand-over & $+$0.9 & $+$0.1 & $-$0.9 & 1.2 & 4/4 \\
rods typed $+5$\,cm, not measured & $+$8.2 & $+$8.2$^b$ & 0.0 & 1.1 & 4/4 \\
attachment radius typed $+2$\,cm & $+$0.9 & 0.0 & $-$0.9 & 1.2 & 4/4 \\
\bottomrule
\end{tabularx}

\end{table}

\paragraph{Rod length and attachment radius.}
\redo{hover with a wrong typed geometry, with and without the rod measurement at the hand-over, for the cases Table~\ref{tab:robust} does not have: rods typed short, the attachment radius typed 2\,cm short, and the combination (rods and radius wrong with the payload mass typed 20\,\% low).}
% --- original text, flown before the rig twin; kept to refill this place:
% \paragraph{Rod length and attachment radius.} The tautness gate scales each UAV's cable feedforward
% by its measured rod length over the typed one. A rod 4\,cm shorter than typed therefore reads as
% partly slack, the fleet flies about half of the cable feedforward, and the payload hangs 22\,cm low.
% Measuring the rods at the hand-over (Section~\ref{sec:takeoff}) removes this
% (Table~\ref{tab:rodlen}): rod errors of $\pm 4$\,cm and attachment-radius errors of $\pm 2$\,cm are
% absorbed to within 1\,cm, and with no error the measurement changes nothing. A mass error is not
% addressed by the measurement; with the payload 20\,\% heavier than typed the hover sat 9\,cm lower.
%
% \begin{table}[h]
% \centering
% \caption{Hover height with a wrong typed geometry, with and without measuring the rods
% (software-in-the-loop, target 0.60\,m, two runs each agreeing to 1\,mm).}
% \label{tab:rodlen}
% \begin{tabularx}{\textwidth}{>{\raggedright\arraybackslash}Xccc}
% \toprule
% Error & Measurement off & Measurement on & Measured rod \\
% \midrule
% none & 0.634 & 0.634 & 0.500 \\
% rod 0.46\,m, typed 0.50\,m & 0.416 & 0.631 & 0.460 \\
% rod 0.50\,m, typed 0.46\,m & 0.569 & 0.634 & 0.500 \\
% radius 0.27\,m, typed 0.25\,m & -- & 0.642 & 0.514 \\
% radius 0.23\,m, typed 0.25\,m & -- & 0.626 & 0.486 \\
% payload mass +20\,\% & -- & 0.540 & 0.500 \\
% rod $-4$\,cm, radius $-2$\,cm, mass +20\,\% & -- & 0.549 & 0.474 \\
% \bottomrule
% \end{tabularx}
% \end{table}

\paragraph{Mass and rod length during an attach.}
\redo{a wrong payload mass ($\pm 25$\,\%) or rod length ($\pm 10$\,\%) in the controller during an attach, two runs each: whether the flight completes, and the payload height after the weld.}
% --- original text, flown before the rig twin; kept to refill this place:
% \paragraph{Mass and rod length during an attach.} On the attach demonstration the controller was told
% a wrong mass or rod length while the simulated plant kept the true ones (two runs each). Telling it
% the payload is 25\,\% lighter, or the rods 10\,\% shorter, than they are was absorbed: both runs
% completed. Telling it the payload is 25\,\% heavier over-lifted the payload to 0.96\,m, and telling it
% the rods are 10\,\% longer dropped it to 0.17\,m; both aborted within about 10\,s of the weld. The
% envelope is asymmetric. An underestimate is made up by the height trim, while an overestimate plans
% more pull than the payload needs. On the rig the payload is weighed and the rods measured, and a
% doubtful value is entered low.

% CUT 5 Oct 2026 (legacy-only, not to be redone): the thrust-gain trim (kt_trim), a legacy-plant feature that is off on the rig and the twin.
% \paragraph{Thrust map.} A per-UAV estimate of the throttle-to-acceleration gain, computed from the
% UAV's own measured acceleration and its static share of the cable force, was built to correct a
% typed gain that is wrong. In the software-in-the-loop bench it brought the hover from 0.499 and
% 0.461\,m back to 0.614\,m with the gain typed 10 and 15\,\% high, with the estimates converging to
% within 1\,\% of the true gain, 3--9\,s after TAKEOFF and without overshoot. In Gazebo, from a floor start with the
% gain typed 10\,\% high, the fleet parked at 0.53\,m without it and lifted to 0.609\,m with it. It needed
% a floor-contact test, because while a UAV rests on the floor the floor carries part of its weight
% and the estimate read that as strong motors. With two UAVs it chattered at 2\,Hz until its inputs were
% time-aligned. On the rig the thrust map was instead identified directly (Section~\ref{sec:hardware}),
% and the estimate is kept off.

% CUT 5 Oct 2026 (legacy-only, not to be redone): the first height-integral test, superseded by Table tab:robust and the disturbance table.
% \paragraph{Height integral.} The bounded height integral on the target (Section~\ref{sec:disturbance})
% was first tested on a hover that sat high without it: 0.658\,m in two runs without the integral
% against 0.600\,m in three runs with it. In software-in-the-loop it held 0.600\,m in every static-offset
% case in both directions, stayed inert through the circle, and saturated with a warning when the error
% was a large gain fault rather than an offset.
